\section{Computational study: implementing the robust counterpart}
\label{sec:experiment}

The preceding section was written, in the first version of this report, as a
theory-only contribution. This section removes that restriction: it implements
the robust placement MILP of \cref{eq:bs-milp} and evaluates it under the same
temporal protocol as \cref{sec:empirical}, on iscf480. Three questions drive
the design. First, is the counterpart \emph{implementable at all} under the
project's own capacity convention --- the $\sigma = 1.10$ slack rule that
calibrates $T_s$? Second, when it is implementable, what does the protection
buy out of sample, and at what price in visits? Third, does \emph{targeted}
budgeted protection beat the obvious nominal alternative --- uniformly
tightening the training ceiling by the same amount of visit price?

\subsection{Setup}
\label{sec:exp-setup}

\paragraph{Solver.}
No commercial solver is available on the experiment machine, so
\cref{eq:bs-milp} is implemented on the open-source HiGHS solver
(\texttt{src/milp\_highs\_robust.py}, via \texttt{highspy}). The visit
variables are declared continuous, which is exact: for fixed integral
$x_{ps}$ the minimal $y_{os} = \max_{p \in P_o} x_{ps}$ is automatically
binary, so the model carries only the $|P|\,|S| = 3{,}840$ assignment binaries
plus the continuous dual block of \cref{prop:bs-counterpart}. Exact
branch-and-bound is nonetheless out of reach at useful time scales: the eight
stations of iscf480 are identical, and the resulting symmetry left HiGHS
unable to improve --- or even find --- an incumbent within minutes. Incumbents
therefore come from a three-stage pipeline: (i) a feasibility-first greedy
(best-fit decreasing on the \emph{robust} station load, swap repair); (ii) a
deterministic pairwise-swap local search on the weighted visit objective, with
the \emph{original} robust constraint --- nominal load plus the sum of the
$\Gamma$ largest deviations, per station --- as a hard accept filter (three
seeded restarts, pass-budgeted for cross-process reproducibility); (iii) a
warm-started HiGHS run used for what an exact solver is still good for here:
\emph{proofs of infeasibility} for budgets that do not fit, and dual bounds.
The optimization objective aggregates the top 1{,}500 weighted distinct
multi-item supports of the training orders; all reported visit counts,
however, are evaluated on the \emph{full real} order sets, and all feasibility
metrics on the fixed $1.10$-rule contract
$T^{\mathrm n} = \lceil 1.10 \sum_p L^{te}_p / (V\,|S|) \rceil$ of
\cref{sec:empirical}. Every solved robust layout was re-verified against the
pre-dualization constraint by direct top-$\Gamma$ arithmetic
(\texttt{robust\_ok} in the result tables): 21 of 21 feasible robust arms
pass, so the dual block never ``cheats.''

\paragraph{Deviations: ex-ante M3 with partial insurance.}
The deviations $\hat L_p$ are calibrated with the only model of
\cref{sec:calibration} available at decision time: M3, the per-product
ddof-$1$ standard deviation of the five volume-rescaled temporal train-block
line counts, capped at $\bar L_p$. (The global-$c$ M2 fit degenerates here:
temporally concentrated SKUs push the $\bar L$-weighted quantile to
$c \approx 70$--$97$, at which point the cap binds for essentially the whole
catalogue and the set collapses to the Soyster box.) Because
\cref{sec:gamma} showed M3 \emph{underestimates} realized drift on iscf480,
and because --- as \cref{sec:exp-impl} shows --- even M3-sized deviations do
not fit the slack contract in full, the experiment adds a partial-insurance
factor $\alpha \in \{0.5, 0.25\}$: the robust constraint is built with
$\alpha \hat L_p$ in place of $\hat L_p$. Formally this is just
\cref{eq:bs-milp} on a shrunken uncertainty set, so
\cref{prop:bs-counterpart} applies verbatim; operationally it asks what
protection is \emph{purchasable} inside the existing capacity margin.

\paragraph{Arms and pairing.}
For each of the five expanding temporal cuts of \cref{sec:empirical}:
the nominal arm $\Gamma = 0$; budgeted arms
$\Gamma \in \{1, 2, 4, 8\} \times \alpha \in \{0.5, 0.25\}$; full-$\sigma$
probes $\Gamma \in \{1, 4, 8, 16\}$ at $\alpha = 1$ (cut~0); two
\emph{uniform-tightening} competitors, trained with the ceiling
$T_s(\beta) = \lceil \beta \sum_p \bar L_p / (V |S|) \rceil$ for
$\beta \in \{1.05, 1.02\}$ and $\Gamma = 0$ --- the same visit price paid as
undirected balance instead of directed protection; and the stored Hexaly
$k{=}1$ layout as an external reference. Every arm of a cut shares the same
supports, budgets, calibration and evaluation, so within-cut comparisons are
paired; the price of robustness $\mathrm{PoR} = (V_{tr}(\cdot) -
V_{tr}(\Gamma{=}0))/V_{tr}(\Gamma{=}0)$ and the test-visit gain
$G = (V_{te}(\Gamma{=}0) - V_{te}(\cdot))/V_{te}(\Gamma{=}0)$ are measured
against the paired nominal arm.

\subsection{Full-\texorpdfstring{$\sigma$}{sigma} protection does not fit the slack contract}
\label{sec:exp-impl}

The $1.10$ rule leaves each station a nominal headroom of
$(1 - 1/1.10)\,T_s$ --- about $4{,}124$ lines per station on cut~0
($T_s = 45{,}365$). The M3 deviations are far larger than that: the top SKU
alone carries $\hat L_p = 13{,}307$ (its block-$\sigma$ exceeds its own
nominal, so the cap binds), and the eight largest deviations sum to
$48{,}621$ --- one and a half times the \emph{total} headroom
$8 \times 4{,}124 \approx 33{,}000$ of the whole cell. Since the heavy
deviators are also heavy products, the nominal constraint prevents parking
them all in one station, and full-$\sigma$ protection has, arithmetically,
nowhere to live. HiGHS makes this formal at moderate budgets:
at $\alpha = 1$ it returns a \textbf{proven infeasibility} certificate for
$\Gamma = 4$, $8$ and $16$ on cut~0, within a minute per arm ($28$--$41$\,s
of total arm time including the constructive pipeline). The
$\Gamma = 1$ case is harder --- its LP relaxation is feasible --- and remains
\textbf{undetermined}: no incumbent and no certificate after a $318$-second
probe and a separate $820$-second attempt (solver logs). The same caution applies at
$\alpha < 1$: the arms marked infeasible in \cref{tab:exp-summary}
($\Gamma \ge 4$ at $\alpha = 0.5$, $\Gamma = 8$ at both $\alpha$, and
$\Gamma = 2$ at $\alpha = 0.5$ on four of five cuts) are
\emph{undetermined-within-budget} --- neither the constructive pipeline nor a
$180$-second HiGHS probe produced a layout --- not proven infeasible. At
half-$\sigma$, by contrast, the top-eight deviation mass is $24{,}310 <
33{,}000$, and indeed $\Gamma = 1$, $\alpha = 0.5$ is constructively feasible
on all five cuts.

The practical reading is blunt: under the project's own capacity convention,
the drift that iscf480 actually exhibits (\cref{sec:gamma}) is too large to be
insured in full. Either the protection target is scaled down ($\alpha$), or
the slack $\sigma$ --- a capacity investment --- must grow. The remainder of
the section quantifies the first option.

\subsection{Results}
\label{sec:exp-results}

\begin{table}[htbp]
\centering\small
\caption{Per-arm means over the temporal cuts on which the arm is feasible
(iscf480; evaluation on real test orders against the fixed $1.10$ contract
$T^{\mathrm n}$). PoR and $G$ are paired against the $\Gamma{=}0$ arm of the
same cut; $G < 0$ means the arm pays visits out of sample too. Arms absent
from the body: at $\alpha=1$, $\Gamma \in \{4,8,16\}$ \emph{proven} infeasible
(HiGHS certificates, cut~0) and $\Gamma = 1$ undetermined; at $\alpha=0.5$,
$\Gamma = 4, 8$ undetermined on 5/5 cuts; at $\alpha = 0.25$, $\Gamma = 8$
undetermined on 5/5 cuts.}
\label{tab:exp-summary}
\begin{tabular}{lrrrrrrr}
\toprule
arm & cuts & PoR\,\% & $G$\,\% & viol.\ ($/8$) & peak $W_s/T^{\mathrm n}$ & worst peak & $X$\,\% \\
\midrule
$\Gamma{=}0$ (nominal)            & 5/5 & 0.00  & 0.00   & 2.6 & 1.77 & 2.11 & 17.4 \\
$\Gamma{=}1$, $\alpha{=}0.5$      & 5/5 & 2.34  & $-$4.44  & 2.6 & \textbf{1.59} & 1.87 & \textbf{13.0} \\
$\Gamma{=}1$, $\alpha{=}0.25$     & 5/5 & 1.44  & $-$2.49  & 2.6 & 1.67 & 1.90 & 15.4 \\
$\Gamma{=}2$, $\alpha{=}0.25$     & 5/5 & 1.41  & $-$1.64  & 2.2 & 1.73 & 1.94 & 15.5 \\
$\Gamma{=}4$, $\alpha{=}0.25$     & 5/5 & 1.77  & $-$2.41  & 2.8 & 1.74 & 2.11 & 15.2 \\
$\Gamma{=}2$, $\alpha{=}0.5$      & 1/5 & 19.13 & $-$23.14 & 3.0 & 1.14 & 1.14 & 3.1 \\
tightened $\beta{=}1.05$          & 5/5 & 0.41  & $-$0.28  & 2.8 & 1.76 & 2.03 & 16.8 \\
tightened $\beta{=}1.02$          & 5/5 & 2.26  & $-$1.73  & 2.8 & 1.75 & 2.01 & 16.9 \\
Hexaly $k{=}1$ (reference)        & 5/5 & $-$3.88 & 6.16 & 2.8 & 2.15 & 2.52 & 23.1 \\
\bottomrule
\end{tabular}
\end{table}

\begin{figure}[htbp]
\centering
\includegraphics[width=\textwidth]{fig_exp_tradeoff}
\caption{Trade-off along the budget $\Gamma$ (means over cuts; open markers =
arms feasible on fewer than five cuts, annotated with their coverage). Left:
price of robustness in train visits. Center: stations above the normalized
ceiling. Right: peak load ratio; the dashed line is contract feasibility.
Partial insurance moves the peak, not the violation count.}
\label{fig:exp-tradeoff}
\end{figure}

\begin{figure}[htbp]
\centering
\includegraphics[width=0.72\textwidth]{fig_exp_frontier}
\caption{Price--severity frontier: mean out-of-sample peak ratio versus PoR.
Budgeted arms (vermillion) move down-right; uniform tightening (blue squares)
moves right \emph{without} moving down; the Hexaly reference (green) shows
that a stronger visit optimizer concentrates load harder. The open marker is
the single-cut $\Gamma{=}2$, $\alpha{=}0.5$ arm.}
\label{fig:exp-frontier}
\end{figure}

\begin{figure}[htbp]
\centering
\includegraphics[width=\textwidth]{fig_exp_stations}
\caption{Per-station load ratios $W_s/T^{\mathrm n}$ on temporal cut~4 (the
worst nominal peak), by arm. The budget compresses the top of the
distribution --- the station a drifting SKU overloads --- rather than
shifting all stations uniformly.}
\label{fig:exp-stations}
\end{figure}

\paragraph{(i) Cheap partial insurance works on severity.}
The efficient point of the sweep is $\Gamma = 1$ at $\alpha = 0.5$: for a
$2.34\%$ visit price it cuts the mean out-of-sample peak from $1.77$ to
$1.59$ and the excess workload fraction $X$ from $17.4\%$ to $13.0\%$
(\cref{tab:exp-summary}). The reduction is \emph{paired sign-consistent}: the
peak ratio falls in five cuts out of five ($1.49\!\to\!1.47$,
$1.48\!\to\!1.34$, $1.81\!\to\!1.61$, $1.95\!\to\!1.87$, $2.11\!\to\!1.64$),
as does $X$. Halving the insured magnitude again ($\alpha = 0.25$) roughly
halves both the price and the effect.

\paragraph{(ii) Equal price, uniform tightening buys nothing.}
The $\beta = 1.02$ competitor pays essentially the same price as the
efficient budgeted arm ($2.26\%$ vs.\ $2.34\%$) and leaves the mean peak at
$1.75$ --- indistinguishable from the nominal $1.77$ --- with per-cut changes
that are sign-\emph{inconsistent} (the peak worsens on two of five cuts).
The mechanism is visible in \cref{fig:exp-stations}: tightening rebalances
\emph{nominal} loads, but the out-of-sample peak is created by large
deviators, and a uniformly tighter ceiling does nothing to separate them ---
the budget constraint does exactly that, because co-locating two large
$\hat L_p$ makes the station's top-$\Gamma$ protection unaffordable. On this
instance family, at equal visit price, \textbf{targeted budgeted protection
dominates uniform tightening on peak severity}.

\paragraph{(iii) The feasibility wall is sharp.}
One step further, $\Gamma = 2$ at $\alpha = 0.5$, is constructible on a
single cut --- and there it nearly restores contract feasibility (peak
$1.14$, $X = 3.1\%$) at a $19.1\%$ visit price. Insurance beyond the slack
headroom is not gradually more expensive; it is discontinuously so, which is
the operational face of the arithmetic in \cref{sec:exp-impl}.

\paragraph{(iv) Severity, not incidence.}
Violation \emph{counts} barely move across implementable arms
($2.2$--$2.8$ of $8$; the count can even rise while the peak falls, as on
cut~3). This is exactly what \cref{sec:gamma} predicts: M3 underestimates the
realized drift, so ex-ante-calibrated protection trims how far stations
overshoot, not whether some station somewhere overshoots. Restoring
feasibility outright would need the drift-scale magnitudes
($c_1(0.9) \approx 1.2$--$1.6$) that no train-internal model provides.

\paragraph{(v) The visit--fragility trade-off, priced.}
The Hexaly reference beats our incumbents by $3.88\%$ on train visits (and
$6.16\%$ on test) --- and carries the worst workload profile of the table
(peak $2.15$, $X = 23.1\%$). Better visit optimization concentrates
correlated, drifting SKUs; the budget constraint is the first dial in this
project that prices that concentration explicitly instead of discovering it
out of sample.

\subsection{Solver honesty}
\label{sec:exp-honesty}

Three caveats bound the claims. First, incumbents are metaheuristic: on the
truncated support objective of cut~0 the HiGHS dual bound is $30{,}312$
against an incumbent of $43{,}158$ ($\approx 30\%$ gap, largely inherited
from the weak LP relaxation of the visit objective), so PoR magnitudes carry
solver noise --- the paired, sign-consistent within-cut deltas are the
reliable signal, not third digits. Second, HiGHS wall-times overran its
\texttt{time\_limit} option on this model (up to $820$\,s against a
$480$\,s limit in the solver logs; persisted arm times reach $465$\,s
against $180$\,s probes), which is why the sweep runs on the deterministic
local-search pipeline and reserves HiGHS for certificates and bounds. Third,
$G < 0$ throughout: every robust arm pays its visit price out of sample as
well ($-1.6\%$ to $-4.4\%$ at the implementable budgets). Nothing here claims
a visits benefit --- the insured quantity is workload-feasibility risk, in
line with the falsified-H1 lesson that opened this report.
